\section{扩展：标注数据重训练（非主贡献）}

本节仅作为扩展讨论，不进入 RQ1--RQ3 的主结论。

若开展重训练，目标是利用反例案例库筛选更有信息量的数据，减少结构性错误（如跨门扩张与配对失败）并提高可计算性/可渲染性。门控失败率的变化仅作为安全性/数据可用性审计与诊断指标披露，不作为主收益主张来源。评估应固定测试集以避免数据泄漏，对比旧模型 $M_0$ 与新模型 $M_1$，并报告 \texttt{model\_issue} 频率、$\mathrm{IoU}_{\mathrm{edit}}$ 的变化，以及安全性事件（门控失败/字段缺失等）的比例与原因分布。

本文当前不做重训练的原因在于样本量与周期有限，且重训练会显著增加方案自由度与泄漏风险控制成本，削弱本文对规程、审计与路由证据链的聚焦。后续若扩展到更大规模且跨数据源的数据，可在保持与本文一致的预注册与隔离协议下再引入重训练闭环。


